\section{Principles}

The security of \texttt{fsb} rests on a few principles that the rest of the design
applies.  Each was, at some point, the lesson of a defect (Section~\ref{sec:history}).

\begin{description}
  \item[One chokepoint.]  Every read of the user's file system goes through one
        package, the \emph{guard}.  No package that serves HTTP may touch the file
        system, and a test fails the build if one imports \code{os}, \code{io/fs},
        \code{syscall} or \code{path/filepath}.  A property of the guard therefore
        holds for every endpoint at once.
  \item[Fail closed.]  When a decision cannot be made, the answer is refusal:
        a malformed rule file stops the program from starting; a file with several
        names is refused until its names are known; a path that cannot be resolved is
        not served.
  \item[Over-match to deny, decide by identity to allow.]  Comparing names can be
        wrong in either direction.  A deny check may safely match too much (it only
        blocks more); an allow check must not, and so decides by the identity of the
        file or folder (device and inode), never by the spelling of its name.
  \item[Name equality belongs to the volume.]  Whether two names are the same file is
        a fact about the file system, not about the program.  The program compares
        names the way the volume does, and tests that claim against the operating
        system rather than against its own model.
  \item[Decide from bytes, not names.]  What a file \emph{is} (an image, a PDF, an
        archive) is decided from its contents.  A name is at most a hint.
  \item[Read nothing that is not needed.]  Refusal is decided before any byte is
        read; a cloud-only file is never read implicitly, because reading it downloads
        it.
  \item[Each layer assumes the one before it failed.]  Rendered Markdown, for
        example, is escaped, then sanitized, then shown in a frame that cannot run
        anything of its own or reach the network.
  \item[Verify against the environment.]  The tests that matter most let the
        operating system decide the facts (Section~\ref{sec:verification}).
\end{description}

\subsection{Technology}

The back end is Go, using the standard library's HTTP server, with one extra
dependency for Unicode (\code{golang.org/x/text}, for normalization and case
folding) and one for extended attributes (\code{golang.org/x/sys}).  The front end
is Svelte 5 with TypeScript, built with Vite, and embedded in the binary with
\code{go:embed}; the compiled output is committed so that building the binary needs
no Node.  The front end has three runtime dependencies that ship in the output:
\code{marked} (Markdown), \code{dompurify} and \code{highlight.js}; Svelte is compiled away.
The result is a single binary of about 11~MB.
